% =============================================================================
% test-abnt-indice.tex
% Documento de teste para o pacote abnt-indice.sty
%
% Demonstra todos os recursos do pacote conforme NBR 6034:2004.
% Cada bloco de texto indica qual requisito da norma está sendo exercitado.
%
% Compilação:
%   Requer shell-escape para que imakeidx chame makeindex automaticamente.
%   TEXINPUTS=./src: latexmk -pdf -pdflatex="pdflatex -shell-escape" -output-directory=tests tests/test-abnt-indice.tex
% =============================================================================

\documentclass[a4paper, 12pt]{report}
\usepackage[T1]{fontenc}
\usepackage[brazil]{babel}
\usepackage[titulo={ÍNDICE DE ASSUNTOS}]{abnt-indice}

\begin{document}

% =====================================================================
% CAPÍTULO 1 — Entradas com acentos e subcabeçalhos
% Exercita: ordenação de acentos (checklist 1), subcabeçalhos (2),
%           cabeçalhos de grupo (8)
% =====================================================================
\chapter{Fundamentos da educação}

A educação\index{educação} no Brasil passou por diversas transformações ao longo
dos séculos. Desde o período colonial, quando as primeiras escolas\index{escola}
foram estabelecidas pelos jesuítas, até as reformas contemporâneas, o sistema
educacional reflete as tensões sociais e políticas do país.

O ensino superior\index{educação!ensino superior} expandiu-se significativamente
a partir da década de 1960, com a criação de novas universidades\index{universidade}
públicas e privadas. Já o ensino fundamental\index{educação!ensino fundamental}
passou por reformas que ampliaram sua duração de oito para nove anos.

A alfabetização\index{alfabetização} é um dos desafios históricos. Diversas
metodologias\index{metodologia} de ensino foram experimentadas, incluindo o método
Paulo Freire\index{Freire, Paulo}, que propõe uma pedagogia\index{pedagogia} baseada
na conscientização do educando.

As áreas\index{área de conhecimento} de conhecimento foram organizadas pela CAPES em
grandes áreas: Ciências Exatas\index{área de conhecimento!Ciências Exatas},
Ciências Humanas\index{área de conhecimento!Ciências Humanas} e
Ciências Biológicas\index{área de conhecimento!Ciências Biológicas}, entre outras.

% =====================================================================
% CAPÍTULO 2 — Remissivas "ver" e "ver também", qualificadores
% Exercita: remissiva "ver" (3), remissiva "ver também" (4),
%           qualificador entre parênteses (7)
% =====================================================================
\chapter{Políticas públicas}

As políticas\index{política educacional} educacionais no Brasil são definidas pelo
Ministério da Educação\index{MEC|see{Ministério da Educação}} (MEC). O
Ministério da Educação\index{Ministério da Educação} atua em conjunto com o Conselho
Nacional de Educação\index{Conselho Nacional de Educação} na formulação de diretrizes.

A avaliação\index{avaliação} institucional tornou-se obrigatória com a criação do
SINAES\index{SINAES|see{Sistema Nacional de Avaliação}}. O
Sistema Nacional de Avaliação\index{Sistema Nacional de Avaliação} foi instituído
pela Lei 10.861/2004.

% Remissiva "ver também" — seção 6.11
\indiceVerTambem{avaliação}{qualidade do ensino}

O Plano Nacional de Educação\index{Plano Nacional de Educação} (PNE) estabelece metas
decenais para a educação. As metas incluem a universalização\index{universalização}
do ensino fundamental e a ampliação do acesso ao ensino superior.

% Qualificador entre parênteses — seção 6.9.2
Pedro II\index{Pedro II (Imperador)} foi responsável pela criação do Colégio
Pedro II\index{Pedro II (Colégio)}, uma das instituições de ensino mais tradicionais
do país, que ainda funciona na cidade do Rio de Janeiro\index{Rio de Janeiro}.

A formação\index{formação de professores} de professores é regulamentada pela
Lei de Diretrizes e Bases\index{Lei de Diretrizes e Bases} da Educação Nacional
(LDB\index{LDB|see{Lei de Diretrizes e Bases}}).

% =====================================================================
% CAPÍTULO 3 — Intervalo de páginas e formatação
% Exercita: páginas consecutivas com hífen (5),
%           páginas não consecutivas com vírgula (6),
%           formatação de página em negrito (9)
% =====================================================================
\chapter{Metodologias de pesquisa em educação}

A pesquisa\index{pesquisa educacional} em educação utiliza tanto métodos
quantitativos\index{metodologia!quantitativa} quanto
qualitativos\index{metodologia!qualitativa}. Os estudos
etnográficos\index{etnografia|textbf} constituem uma das abordagens mais
utilizadas na pesquisa educacional brasileira. A etnografia permite ao
pesquisador compreender as práticas educativas no contexto em que ocorrem.

% Mesmo termo em páginas consecutivas para gerar intervalo — seção 6.12a
A análise\index{currículo|(} do currículo é tema central nas pesquisas sobre
educação. O currículo escolar define os conteúdos e as competências a serem
desenvolvidas pelos estudantes.

\newpage

As reformas curriculares acompanham as mudanças nas demandas sociais. A Base
Nacional Comum Curricular\index{BNCC|see{currículo}} (BNCC) é o documento mais
recente que orienta a organização do currículo\index{currículo|)} nas escolas
brasileiras.

A inclusão\index{inclusão} escolar é outro tema de pesquisa relevante. As
políticas de inclusão visam garantir o acesso de estudantes com
deficiência\index{deficiência} ao ensino regular.

% =====================================================================
% CAPÍTULO 4 — Mais entradas para variedade de grupos alfabéticos
% Exercita: cabeçalhos de grupo (8), título descritivo (10)
% =====================================================================
\chapter{Perspectivas contemporâneas}

A tecnologia\index{tecnologia educacional} tem transformado as práticas pedagógicas.
O uso de plataformas\index{plataforma digital} digitais ampliou o acesso à
educação a distância\index{educação!a distância}, especialmente durante o período
de isolamento\index{isolamento social} social.

A qualidade do ensino\index{qualidade do ensino} é medida por indicadores como o
IDEB\index{IDEB|see{qualidade do ensino}} (Índice de Desenvolvimento da Educação
Básica). Os resultados\index{resultados educacionais} mostram desigualdades
regionais significativas.

A gestão\index{gestão escolar} escolar democrática é um princípio estabelecido pela
Constituição\index{Constituição Federal} de 1988. Os conselhos
escolares\index{conselho escolar} representam um dos mecanismos de participação da
comunidade nas decisões da escola.

A valorização\index{valorização docente} dos profissionais da educação inclui
questões de remuneração\index{remuneração}, formação continuada\index{formação de
professores} e condições de trabalho. O Piso Salarial Nacional\index{Piso Salarial
Nacional} foi estabelecido pela Lei 11.738/2008.

% Referência não consecutiva — seção 6.12b
% "universidade" aparece no cap. 1 (p.1) e aqui (p.4+), gerando vírgula
As universidades\index{universidade} desempenham papel fundamental na produção de
conhecimento e na formação de pesquisadores\index{pesquisador}.

% =====================================================================
% ÍNDICE
% =====================================================================

\imprimirindice

\end{document}
